\begin{name}
	{\tenchude}
	{TOÁN 10}
	{LỚP TOÁN THẦY PHÁT}
	{{Ai cũng sẽ phải trưởng thành dù muốn hay không.\\ Vậy hãy chọn cách mạnh mẽ mà trưởng thành.}}
\end{name}
\Opensolutionfile{ansbook}[ans/ansbookDe6]
\TN
\Opensolutionfile{ans}[ans/ansDe1-TN6]
\begin{ex}%[HSA - đợt 2, Nguyễn Thành Nhân]%[0D8N1-2]
Bình có $5$ cái áo khác nhau, $4$ cái quần khác nhau, $3$ đôi giày khác nhau và $2$ chiếc mũ khác nhau. Số cách chọn một bộ gồm một quần, một áo, một đôi giày và một chiếc mũ là
\choice
{\True $120$}
{$5$}
{$60$}
{$14$}
\loigiai{
Theo quy tắc nhân, ta có số cách chọn một bộ là $5\times 4\times 3\times 2=120$.
}
\end{ex}

\begin{ex}%[24-25 giảng K10 - K11, Phan Quốc Trí]%[0D8N2-5]
Có bao nhiêu cách cắm $ 3 $ bông hoa giống nhau vào $ 5 $ lọ khác nhau (mỗi lọ cắm  một bông)?
\choice
{\True $ 10 $}
{$ 30 $}
{$ 6 $}
{$ 60 $}
\loigiai{
Cắm $ 3 $ bông hoa giống nhau, mỗi bông vào $ 1 $ lọ nên ta sẽ lấy $ 3 $ lọ bất kỳ trong $ 5 $ lọ khác nhau để cắm bông. Vậy số cách cắm bông chính là tổ hợp chập $ 3 $ của $ 5 $ phần tử. Như vậy, ta có $ \mathrm{C}_5^3=10 $ cách.
}
\end{ex}

\begin{ex}%[0D8N3-2]%[HKII-NGUYỄN THÁI BÌNH 2324]%[TheHung Nguyen]
Có bao nhiêu số hạng trong khai triển nhị thức $\left(2 x^2+1\right)^4$?
\choice
{$3$}
{$8$}
{$4$}
{\True $5$}
\loigiai{
Khai triển nhị thức $\left(2 x^2+1\right)^4$ có $4+1=5$ số hạng.
}
\end{ex}

\begin{ex}%[H]%[Nguyễn Văn Cường (Cường NV) - Dự án giảng K10 - K11]%[0D6H1-1]
Một hình chữ nhật có các cạnh $x = 4{,}2$ m $\pm 1$ cm, $y = 7$ m $\pm 2$ cm. Chu vi của hình chữ nhật và sai số tuyệt đối của giá trị đó.
\choice
{$22{,}4$ m và $3$ cm}
{$22{,}4$ m và $1$ cm}
{$22{,}4$ m và $2$ cm}
{\True $22{,}4$m và $6$ cm}
\loigiai{
Ta có chu vi hình chữ nhật là $P=2(x+y)=22{,}4$ m $\pm 6$ cm.}
\end{ex}

\begin{ex}%[0-HK1-CT-1-BuiThiXuan-TPHCM-2324]%[VN-MT-6, VM013]%[0D6H4-2]
Số bàn thắng mà một đội bóng ghi được ở mỗi trận đấu trong một mùa giải được thống kê lại ở bảng sau
\begin{center}
\begin{tabular}{|c|c|c|c|c|c|c|}
\hline
Số bàn thắng&$0$&$1$&$2$&$3$&$4$&$5$\\
\hline
Số trận&$4$&$8$&$6$&$3$&$2$&$1$\\
\hline
\end{tabular}
\end{center}
Tìm khoảng tứ phân vị của số bàn thắng trong bảng trên
\choice
{$1$}
{\True $1{,}5$}
{$3$}
{$2$}
\loigiai{
Cỡ mẫu của mẫu số liệu trên là $n=24$.\\
Trung vị (tứ phân vị thứ hai) của mẫu số liệu trên là \[Q_2=\dfrac{x_{12}+x_{13}}{2}=\dfrac{1+2}{2}=1{,}5.\]
Trung vị cho nửa mẫu số liệu bên trái (tứ phân vị thứ nhất) là \[Q_1=\dfrac{x_6+x_7}{2}=\dfrac{1+1}{2}=1.\]
Trung vị cho nửa mẫu số liệu bên phải (tứ phân vị thứ ba) là \[Q_3=\dfrac{x_{18}+x_{19}}{2}=\dfrac{2+3}{2}=2{,}5.\]
Khoảng tứ phân vị của mẫu số liệu là \[\triangle{Q}=Q_3-Q_1=2{,}5-1=1{,}5.\]
}
\end{ex}

\begin{ex}%[0D0N1-2]%[Dự án đề kiểm tra Toán khối 10 HKII NH23-24-Đợt 15-Lê Hải Phụng]%[THPT Chuyên Lê Quý Đôn - HCM]
Gieo một con xúc xắc cân đối và đồng chất hai lần. Xác định số phần tử của không gian mẫu.
\choice
{$ n\left(\Omega\right)=6$}
{\True $ n\left(\Omega\right)=36$}
{$ n\left(\Omega\right)=2$}
{$ n\left(\Omega\right)=12$}
\loigiai{Gieo một con xúc xắc cân đối và đồng chất hai lần thì $ n\left(\Omega\right)=6\cdot6 =36$.}
\end{ex}

\begin{ex}%[0D0N1-3]%[Dự án đề kiểm tra Toán 10 GHKI NH23-24- Lam Nguyen]%[THPT Le Quý Đôn - TPHCM]
Xét phép thử  \lq\lq Gieo một xúc xắc hai lần liên tiếp \rq\rq. Gọi $A$ là biến cố \lq\lq Tổng số chấm hai lần gieo không nhỏ hơn $10$ \rq\rq. Số kết quả thuận lợi cho $A$ là
\choice
{$3$}
{\True $4$}
{$5$}
{$6$}
\loigiai{Số kết quả thuận lợi của biến cố $A$ là các phần tử $(5;5)$, $(6;5)$, $(5;6)$, $(6;6)$.
}
\end{ex}

\begin{ex}%[10-11-12EX-GHK2-2425]%[Dương Quang]%[0H9H3-1]
Trong mặt phẳng $Oxy$, cho đường thẳng $d\colon\heva{&x=2-t \\&y=5+3t}$ ($t\in\mathbb{R}$). Điểm nào sau đây thuộc đường thẳng?
\choice
{$P(2;8)$}
{\True $Q(1;8)$}
{$N(1;5)$}
{$Q(2;-5)$}
\loigiai{
Tại $t=1$ thì $\heva{&x=1 \\&y=8}$. Do đó $Q(1;8)\in d$.
}
\end{ex}

\begin{ex}%[0H9H3-2]%[Dự án đề kiểm tra Toán khối 10 HK2-NH23-24-Đợt 15-Dương Công Tạo]%[THPT Đặng Huy Trứ - Huế]
Cho đường thẳng $d$ có phương trình tham số $\heva{&x=5+2t \\&y=-1-t.&}$ Phương trình tổng quát của đường thẳng qua $B(1; 1)$ và vuông góc với đường thẳng $d$ là
\choice
{$2x-y-3=0$}
{$x+2y-3=0$}
{\True $2x-y-1=0$}
{$2x+y-3=0$}
\loigiai{
Gọi $\Delta$ là đường thẳng cần tìm.\\
Ta có $\Delta \perp d$ suy ra véc-tơ pháp tuyến của $\Delta$ chính là véc-tơ chỉ phương của $d$.\\
Tức là $\vec{n}_{\Delta}=\vec{u}_d=(2;-1)$.\\
Khi đó, phương trình tổng quát của đường thẳng $\Delta$ đi qua $B(1; 1)$ và nhận $\vec{n}_{\Delta}=(2;-1)$ làm véc-tơ pháp tuyến là \[2(x-1)-1(y-1)=0\text{ hay } 2x-y-1=0.\]
}
\end{ex}

\begin{ex}%[0H9H3-4]
Cho đường thẳng $d_{1} \colon 3x+4y+1=0$ và $d_{2}  \colon \heva{& x=15+12t \\& y=1+5t}$. Tính cosin của góc tạo bởi giữa hai đường thẳng đã cho.
\choice{$\dfrac{56}{65}$}
{$-\dfrac{33}{65}$}
{$\dfrac{6}{65}$}
{\True $\dfrac{33}{65}$}
\loigiai{
Ta có $\overrightarrow{n_{1}}=\left(3;4\right) \Rightarrow \overrightarrow{u_{1}}=\left(-4;3\right); \overrightarrow{u_{2}}=\left(12;5\right)$.\\
$\cos \left(d_{1} ;d_{2} \right)=\left|\cos \left(\overrightarrow{u_{1}};\overrightarrow{u_{2}}\right)\right|=\dfrac{\left|-4.12+3.5\right|}{\sqrt{\left(-4\right)^{2} +3^{2}} .\sqrt{12^{2} +5^{2}}} =\dfrac{33}{65}$.}
\end{ex}

\begin{ex}%[0H9N5-1]%[Dự án đề kiểm tra Toán 10 HKII NH23-24- Đoàn Thanh Phong]%[THPT Thiệu Hóa - Thanh Hóa]
Đường Elip $\dfrac{x^2}{16}+\dfrac{y^2}{7}=1$ có tiêu cự bằng
\choice
{$\dfrac{9}{16}$}
{$3$}
{\True $6$}
{$\dfrac{6}{7}$}
\loigiai{
Ta có $a^2=16$, $b^2=7$ nên $c^2=a^2-b^2=9\Rightarrow c=3$.\\
Vậy Elip có tiêu cự bằng $2c=2\cdot3=6$.
}
\end{ex}

\begin{ex}%[0H9N5-5]
Phương trình nào sau đây là phương trình của đường hypebol?
\choice
{\True $\dfrac{x^2}{9} - \dfrac{y^2}{16} = 1$}
{$\dfrac{x^2}{9} + \dfrac{y^2}{16} = 1$}
{$\dfrac{x^2}{9} - \dfrac{y^2}{16} = 0$}
{$\dfrac{x^2}{9} + \dfrac{y^2}{16} = 0$}
\loigiai{
$\dfrac{x^2}{9} - \dfrac{y^2}{16} = 1$ là dạng đúng phương trình của một hypebol.
}
\end{ex}
\TNTF
\begin{ex}%[Mức độ 3]%[0D8V2-5]
Mỗi máy tính tham gia vào mạng phải có một địa chỉ duy nhất, gọi là địa chỉ IP, nhằm định danh máy tính đó trên Internet. Xét tập hợp $A$ gồm các địa chỉ IP có dạng $192.168.abc.deg$, trong đó $a$, $b$, $c$ là các chữ số phân biệt được chọn ra từ các chữ số $0$, $1$, $2$, $3$, $4$, còn $d$, $e$, $g$ là các chữ số phân biệt được chọn ra từ các chữ số $5$, $6$, $7$, $8$, $9$.
\choiceTF[t]
{\True Địa chỉ IP $192.168.123.789$ khác với địa chỉ IP $192.168.321.987$}
{\True Trong địa chỉ IP dạng $192.168.abc.deg$ thì các chữ số ở các vị trí $a$, $b$, $c$, $d$, $e$, $g$ phải hoàn toàn khác nhau}
{Trong địa chỉ Ip dạng $192.168.abc.deg$ thì chữ số $0$ không được đứng ở vị trí $a$}
{Tập hợp $A$ có tất cả $120$ phần tử}
\loigiai
{\begin{itemchoice}
\itemch {\bfseries Đúng.}\\
Địa chỉ IP $192.168.123.789$ khác với địa chỉ IP $192.168.321.987$.
\itemch {\bfseries Đúng.}\\
Theo mô tả, $a$, $b$, $c$ là các chữ số phân biệt được chọn ra từ các chữ số $0$, $1$, $2$, $3$, $4$, còn $d$, $e$, $g$ là các chữ số phân biệt được chọn ra từ các chữ số $5$, $6$, $7$, $8$, $9$ nên các chữ số ở các vị trí $a$, $b$, $c$, $d$, $e$, $g$ phải hoàn toàn khác nhau.
\itemch {\bfseries Sai.}\\
Trong địa chỉ Ip dạng $192.168.abc.deg$ thì chữ số $0$ vẫn có thể đứng ở vị trí $a$.
\itemch {\bfseries Sai.}\\
Mỗi cách chọn chữ số cho $3$ vị trí $a$, $b$, $c$ từ $5$ chữ số $0$, $1$, $2$, $3$, $4$ là một chỉn hợp chập $3$ của $5$ phần tử. Do đó có $\mathrm{A}_5^3=60$ cách.\\
Tương tự, có $\mathrm{A}_5^3=60$ cách chọn chữ số cho $3$ vị trí $d$, $e$, $g$.\\
Như thế, có tất cả $60\times 60=3\,600$ địa chỉ IP.\\
Vậy tập hợp $A$ có tất cả $3\,600$ phần tử.
\end{itemchoice}

}
\end{ex}

\begin{ex}%[0H9H3-2]  %[Dự án đề kiểm tra Toán 10 HKII NH23-24- Thầy Hải Toán]%[THPT Đống Đa Hà Nội]
Cho hai điểm $A(3;-3)$, $B(-1;-5)$ và đường thẳng $(d)\colon 4x-3y-2=0$
\choiceTF
{Đường thẳng đi qua điểm $A$ và vuông góc với $(d)$ có phương trình $4x+3y=3$}
{Đường tròn đường kính $AB$ có phương trình $(x-1)^2+(y+4)^2=20$}
{Khoảng cách từ $A$ tới $(d)$ nhỏ hơn khoảng cách từ $B$ tới $(d)$}
{\True Cosin của góc tạo bởi $(d)$ và đường thẳng $AB$ bằng $\dfrac{2}{\sqrt{5}}$}
\loigiai{
\begin{enumerate}[a)]
\item Đường thẳng vuông góc với đường thẳng $(d)$ nên có dạng: $\Delta\colon 3x+4y+m=0$.\\
$\Delta$ đi qua $A$ nên thay toạ độ $A$ vào phương trình đường thẳng $\Delta$ ta có \[3\cdot3+4\cdot(-3)+m=0\Leftrightarrow m=3.\]
Vậy phương trình đường thẳng cần tìm là $\Delta\colon 3x+4y+3=0$.
\item Đường tròn đường kính $AB$ thì $AB$ đi qua tâm $I$ và $I$ là trung điểm của $AB$ suy ra $I\left(2;-4\right)$. Suy ra $\vv{IA}=\left(1;1\right)$\\
Bán kính $R=\big|\vv{IA}\big|=\sqrt{1^2+1^2}=\sqrt{2}$.\\
Vậy phương trình đường tròn đường kính $AB$ là $(C)\colon \left(x-2\right)^2+\left(y+4\right)^2=2$.
\item Khoảng cách từ $A$ đến đường thẳng $(d)$ là  $\mathrm{d}\left(A,d\right)=\dfrac{\big|4\cdot 3-3\cdot(-3)-2\big|}{\sqrt{4^2+(-3)^2}}=\dfrac{19}{5}$.\\
Khoảng cách từ $B$ đến đường thẳng $(d)$ là $\mathrm{d}\left(B,d\right)=\dfrac{\big|4\cdot(-1)-3\cdot(-5)-2\big|}{\sqrt{4^2+(-3)^2}}=\dfrac{9}{5}$.\\
Vậy $\mathrm{d}\left(A,d\right)>\mathrm{d}\left(B,d\right)$.
\item Đường thẳng $AB$ nhận $\vv{AB}=\left(-4;-2\right)$ làm một véc-tơ chỉ phương.\\
đường thẳng $(d)$ có véc-tơ pháp tuyến là $n_d=\left(4;-3\right)$ suy ra véc-tơ chỉ phương là $u_d=\left(3;4\right)$.\\
Vậy $\cos\left(d,AB\right)=\dfrac{\big|-4\cdot 3+(-2)\cdot4\big|}{\sqrt{\left(-4\right)^2+\left(-2\right)^2}\cdot\sqrt{3^2+4^2}}=\dfrac{2\sqrt{5}}{5}$.
\end{enumerate}
}
\end{ex}
\TNSA
\begin{ex} %[0D8V1-3]
Một bệnh viện có đội ngũ các nhân viên y tế trẻ tham gia đội ngũ tình nguyện viên gồm $12$ bác sĩ nội khoa, $10$ bác sĩ ngoại khoa. Có bao nhiêu cách chọn một đội gồm $5$ người trong đó phải có bác sĩ nội khoa, bác sĩ ngoại khoa?
\loigiai{
Việc chọn một đội bác sĩ tình nguyện có các trường hợp sau đây:
\begin{itemize}
\item Có $\mathrm{C}_{12}^4\cdot\mathrm{C}_{10}^1=4\,950$ cách chọn đội gồm có $4$ bác sĩ nội khoa, $1$ bác sĩ ngoại khoa.
\item Có $\mathrm{C}_{12}^3\cdot\mathrm{C}_{10}^2=9\,900$ cách chọn đội gồm có $3$ bác sĩ nội khoa, $2$ bác sĩ ngoại khoa.
\item Có $\mathrm{C}_{12}^2\cdot\mathrm{C}_{10}^3=7\,920$ cách chọn đội gồm có $2$ bác sĩ nội khoa, $2$ bác sĩ ngoại khoa.
\item Có $\mathrm{C}_{12}^1\cdot\mathrm{C}_{10}^4=2\,520$ cách chọn đội gồm có $1$ bác sĩ nội khoa, $4$ bác sĩ ngoại khoa.
\end{itemize}
Vậy có tất cả $25290$ cách chọn thỏa yêu cầu.
}
\end{ex}

\begin{ex}%[0D8V3-4]
Xác định hệ số của $x^3$ trong khai triển biểu thức $\left(\dfrac{2}{3}x+\dfrac{1}{4}\right)^5$
\loigiai{
Ta có $\left(\dfrac{2}{3}x+\dfrac{1}{4}\right)^5
=\sum\limits_{k=0}^5 \mathrm{C}_5^k\left(\dfrac{2}{3}x\right)^{5-k}\left(\dfrac{1}{4}\right)^k
=\sum\limits_{k=0}^5 \left(\dfrac{2}{3}\right)^{5-k}\left(\dfrac{1}{4}\right)^k\mathrm{C}_5^kx^{5-k}$. \\
Hệ số của $x^3$ ứng với $5-k=3\Rightarrow k=2\Rightarrow a_3=\left(\dfrac{2}{3}\right)^3\left(\dfrac{1}{4}\right)^2\mathrm{C}_5^2=\dfrac{5}{27}$.
}
\end{ex}

\begin{ex}%[H]%[BG10-3in1, Phạm Ngọc Trung]%[0D6H3-2]
Người ta điều tra ngẫu nhiên số cân nặng (đơn vị kg) của $20$ học sinh nữ một trường phổ thông, kết quả được ghi trong bảng sau
\begin{center}
\begin{tabular}{|c|c|c|c|c|c|c|c|}
\hline Số cân nặng & $38$ & $40$ & $43$ & $45$ & $48$ & $50$ & \\
\hline Tần số & $3$ & $3$ & $4$ & $5$ & $2$ & $3$ & $N=20$ \\
\hline
\end{tabular}
\end{center}
Tính số cân nặng trung bình $\overline{x}$ (\textit{kết quả làm tròn đến hàng phần chục}).
\shortans{43{,}9}
\loigiai{Số cân nặng trung bình của $20$ học sinh là
\begin{align*}
\overline{x}=\dfrac{38\cdot 3+40\cdot 3+43\cdot 4+45\cdot 5+48\cdot 2+50\cdot 3}{20}=\dfrac{877}{20}\approx 43{,}9.
\end{align*}
}
\end{ex}

\begin{ex}%[dự án HSA, Huỳnh Xuân Tín]%[0H9V4-2]
Đường tròn $(C)$ có tâm $I$ thuộc đường thẳng $d\colon x+3y+8=0$, đi qua điểm $A\left(-2;1\right)$ và tiếp xúc với đường thẳng $\Delta:3x-4y+10=0$. Tính bán kính đường tròn $(C)$.
\shortans[1]{$5$}
\loigiai{
Ta có $I\in d \Rightarrow I(-3m-8;m)$.
\begin{eqnarray*}
&&AI^2=\left[\mathrm{d}(I,\Delta)\right]^2 \text{ (cùng bằng } R^2) \\
& \Leftrightarrow & \left(-3m-6\right)^2+\left(m-1\right)^2=\dfrac{\left[3(-3m-8)-4m+10\right]^2}{3^2+(-4)^2}\\
& \Leftrightarrow & 10m^2+34m+37=\dfrac{169m^2+364m+196}{25}\\
& \Leftrightarrow & 81m^2+486m+729=0\\
& \Leftrightarrow & m=-3.
\end{eqnarray*}
Do đó, đường tròn $(C)$ có $\heva{&\text{tâm } I(1;-3)\\& \text{bán kính } R=AI=\sqrt{(9-6)^2+(-3-1)^2}=5.}$\\
Vậy phương trình đường tròn là ${\left(x-1\right)}^2+{\left(y+3\right)}^2=25$.
}
\end{ex}
\TL
\begin{ex}%[0H9V3-2]%[Dự án đề kiểm tra Toán 10 CHKII NH23-24- Ngô Tất Thành]%[THPT Chuyên Hùng Vương - Tỉnh Phú Thọ]
Lập phương trình đường thẳng $\Delta$ đi qua điểm $N(5;2)$ và song song với đường thẳng $3x-2y+5=0$.
\loigiai{
Do đường thẳng $\Delta$ song song với đường thẳng $3x-2y+5=0$ nên phương trình của $\Delta$ có dạng
\[3x-2y+c=0 \qquad (c \ne 5)\]
Đường thẳng $\Delta$ đi qua điểm $N(5;2)$
\[\Leftrightarrow 3 \cdot 5 -2 \cdot 2 +c=0 \Leftrightarrow c=-11(TM).\]
Vậy phương trình đường thẳng $\Delta$ là $3x-2y-11=0$.
}
\end{ex}

\begin{ex}%[0H9V4-5]
Trong mặt phẳng với hệ trục tọa độ $Oxy$, cho đường tròn \[(C)\colon x^2+y^2+4x-2y-20=0\] và đường thẳng $\Delta\colon 3x+4y-m=0$. Tìm giá trị dương của tham số $m$ để đường thẳng $\Delta$ cắt đường tròn $(C)$ theo một dây dung có độ dài bằng $6$.
\loigiai{
\begin{center}
\begin{tikzpicture}[declare function={r=2.5;ga=40;},>=stealth,font=\scriptsize,scale=0.8]
%\path (0:0) coordinate (I)
%(ga:r) coordinate (A)
%(130:r) coordinate (B)
\coordinate (I) at (0,0);
\coordinate (A) at (ga:r);
\coordinate (B) at (130:r);
\coordinate (J) at ($(A)!1/2!(B)$);
\draw (I) circle (r)  ;
\draw (I)--(A)--(B)--(I)--(J);
\foreach \x/\g in{A/90,B/90,J/90,I/-90}
\fill[black](\x) circle (2pt)
($(\x)+(\g:5mm)$) node{\small $\x$};
\end{tikzpicture}
\end{center}
Đường tròn $(C)$ có tâm $I(-2;1)$ và bán kính $R=\sqrt{(-2)^2+1^2+20}=5$.\\
Giả sử $\Delta$ cắt $(C)$ tại hai điểm $A$ và $B$. Gọi $J$ là trung điểm của $AB$.\\
Khi đó $\heva{&IJ\perp AB\\&JB=\dfrac{1}{2}AB=3.}$\\
Xét tam giác $IJB$ vuông tại $J$ có $IJ=\sqrt{IB^2-JB^2}=4.$\\
Hay $\mathrm{d}(I; \Delta)=4$. Tương đương
\begin{align*}
&\dfrac{|(-2)\cdot3+4\cdot1-m|}{\sqrt{3^+4^2}}=4\\
\Leftrightarrow&\, |m+2|=4\cdot5=20\\
\Leftrightarrow&\, \hoac{&m+2=20\\&m+2=-20}\Leftrightarrow\hoac{&m=18\\&m=-22\,(\text{không thỏa mãn}).}
\end{align*}
Vậy $m=18$.
}

\end{ex}

\begin{ex}%[Dự án C THPTQG 2025]%[Vương Quốc Phong]%[0D0C2-9]
Hai bạn Nga và Nhung chơi trò tung xúc xắc. Mỗi bạn tung $1$ con xúc xắc $3$ lần, ai có tổng số chấm $3$ lần gieo lớn hơn thì thắng. Nga chơi trước và được $14$ chấm. Tính xác suất để Nhung thắng Nga.
\loigiai{
Gọi $A$ là biến cố: \lq\lq Nhung thắng Nga sau ba lần tung\rq\rq.\\
Khi đó $\mathrm{P}(A)$ là xác suất tổng số chấm Nga tung được sau ba lần tung lớn hơn $14$. \\
Ta có $\Omega=\left\{(a_1, a_2, a_3)|a_i \in \{1,2,\ldots, 6\}, i=\overline{1,3}\right\}\Rightarrow n(\Omega)=6^3=216$.\\
$A=\left\{(a_1, a_2, a_3)|a_1+a_2+a_3 \ge 15, a_i \in \{1,2,\ldots, 6\}, i=\overline{1,3}\right\}$.\\
Để đếm số phần tử của $A$, ta chia thành các trường hợp sau:
\begin{itemize}
\item Trường hợp 1: $a_1+a_2+a_3=15$, gồm bộ $(5,5,5)$ và các bộ là hoán vị của $(4;5;6)$ và $(3;6;6)$.\\
Trường hợp này có $1+6+3=10$ (bộ)
\item  Trường hợp 2: $a_1+a_2+a_3=16$, gồm các hoán vị của $(5,5,6)$ và $(4,6,6)$.\\ Trường hợp này có $3+3=6$ (bộ)
\item  Trường hợp 3: $a_1+a_2+a_3=17$, gồm các hoán vị của $(5,6,6)$.\\
Trường hợp này có $3$ (bộ).
\item Trường hợp 4: $a_1+a_2+a_3=18$, gồm $(6,6,6)$.\\
Trường hợp này có $1$ (bộ).
\end{itemize}
Vậy $n(A)=20$. \\
Xác suất cần tìm là $\mathrm{P}(A)=\dfrac{n(A)}{n(\Omega)}=\dfrac{20}{216}=\dfrac{5}{54}$.
}
\end{ex}
\Closesolutionfile{ans}
\Closesolutionfile{ansbook}